% Project-specific blueprint content for QuadraticNumberFields.

\chapter{Blueprint Scope}

This blueprint is organized around the main theorem of the project: the
classification of the ring of integers of the quadratic field
$\mathbf{Q}(\sqrt d)$ for squarefree integers $d \neq 1$.
The formal development works with two kinds of data:
\begin{description}
\item[Ambient field] The quadratic $\mathbf{Q}$-algebra $Q(\sqrt d)$, modeled in
Lean by \verb|Qsqrtd d|.
\item[Integral models] The two orders $\mathbf{Z}[\sqrt d]$ and
$\mathbf{Z}[(1+\sqrt d)/2]$, which are exactly the two candidates for the full
ring of integers.
\item[Classification mechanism] An integral element is first rewritten in
half-integer coordinates $(a' + b' \sqrt d)/2$, and then the divisibility
condition on $a'^2 - d b'^2$ is analyzed modulo $4$.
\end{description}

The blueprint focuses on the classification pipeline and one structural
corollary, namely the Dedekind-domain criterion for $\mathbf{Z}[\sqrt d]$.
Other parts of the repository, such as discriminants, totally real/complex
classification, and the $\mathbf{Z}[\sqrt{-5}]$ examples, are downstream of this
core story.

\chapter{Basic Models}

\begin{definition}[Quadratic field model]\label{def:qsqrtd}%
\lean{Qsqrtd}\leanok
For $d \in \mathbf{Q}$, define the quadratic field model
\[
  Q(\sqrt d) := \mathrm{QuadraticAlgebra}\,\mathbf{Q}\, d\, 0.
\]
This is the ambient field in which all integral-closure statements are proved.
\end{definition}

\begin{definition}[The Zsqrtd order]\label{def:zsqrtd}%
\lean{QuadraticNumberFields.RingOfIntegers.Zsqrtd}\leanok
For $d \in \mathbf{Z}$, define
\[
  \mathbf{Z}[\sqrt d] := \mathrm{QuadraticAlgebra}\,\mathbf{Z}\, d\, 0.
\]
The file \verb|RingOfIntegers/Zsqrtd.lean| also provides the canonical embedding
of this order into $Q(\sqrt d)$.
\end{definition}

\begin{definition}[Half-integer coordinates]\label{def:halfint}%
\lean{QuadraticNumberFields.RingOfIntegers.halfInt, QuadraticNumberFields.RingOfIntegers.Zsqrtd.halfInt}\leanok
For integers $a', b', d$, define the half-integer element
\[
  \frac{a' + b' \sqrt d}{2} \in Q(\sqrt d).
\]
This normal form is the key bridge between integrality and congruence conditions
modulo $4$.
\end{definition}

\begin{definition}[The Zomega order]\label{def:zomega}%
\lean{ZOnePlusSqrtOverTwo}\leanok
For $k \in \mathbf{Z}$, define
\[
  \mathbf{Z}\Bigl[\frac{1 + \sqrt{1+4k}}{2}\Bigr]
\]
as the quadratic algebra over $\mathbf{Z}$ with relation $\omega^2 = \omega + k$.
This is the correct integral model in the branch $d \equiv 1 \pmod 4$.
\end{definition}

\chapter{Integrality Pipeline}

\section{Trace and Norm}

\begin{lemma}[Integral trace]\label{lem:int-trace}%
\lean{QuadraticNumberFields.RingOfIntegers.TraceNorm.exists_int_trace}\leanok
If $x \in Q(\sqrt d)$ is integral over\/ $\mathbf{Z}$, then
\[
  \mathrm{Tr}_{Q(\sqrt d)/\mathbf{Q}}(x) \in \mathbf{Z}.
\]
Formally, the theorem produces an integer $a'$ such that
$a' = \mathrm{trace}(x)$.
\end{lemma}

\begin{lemma}[Integral norm]\label{lem:int-norm}%
\lean{QuadraticNumberFields.RingOfIntegers.TraceNorm.exists_int_norm}\leanok
If $x \in Q(\sqrt d)$ is integral over\/ $\mathbf{Z}$, then
\[
  N_{Q(\sqrt d)/\mathbf{Q}}(x) \in \mathbf{Z}.
\]
Formally, the theorem produces an integer $n'$ such that
$n' = \mathrm{norm}(x)$.
\end{lemma}

\begin{lemma}[Half-integrality of the imaginary coordinate]\label{lem:int-double-im}%
\lean{QuadraticNumberFields.RingOfIntegers.TraceNorm.exists_int_double_im}\leanok
Once the trace and norm of an integral element are known to be integers, the
imaginary coordinate has denominator at most $2$: there is an integer $b'$ such
that
\[
  \operatorname{im}(x) = \frac{b'}{2}.
\]
This is the final ingredient needed to put $x$ into half-integer normal form.
\end{lemma}

\begin{theorem}[Half-integer normal form for integral elements]%
\label{thm:halfint-normal-form}%
\lean{QuadraticNumberFields.RingOfIntegers.exists_halfInt_with_div_four_of_isIntegral}\leanok
\uses{def:halfint,lem:int-trace,lem:int-norm,lem:int-double-im}
If $x \in Q(\sqrt d)$ is integral over\/ $\mathbf{Z}$, then there exist integers
$a', b'$ such that
\[
  x = \frac{a' + b' \sqrt d}{2}
  \quad\text{and}\quad
  4 \mid (a'^2 - d b'^2).
\]
This theorem is the central normal-form statement in the formalization: it
reduces the classification of integral elements to a congruence problem in
$\mathbf{Z}$.
\end{theorem}

\section{The Non-\texorpdfstring{$1 \bmod 4$}{1 mod 4} Branch}

\begin{lemma}[Modulo-$4$ criterion away from $1 \bmod 4$]%
\label{lem:mod4-ne-one}
\lean{QuadraticNumberFields.RingOfIntegers.dvd_four_sub_sq_iff_even_even_of_ne_one_mod_four}\leanok
Assume $d$ is squarefree and $d \not\equiv 1 \pmod 4$. Then
\[
  4 \mid (a'^2 - d b'^2)
  \iff
  2 \mid a' \text{ and } 2 \mid b'.
\]
This is the arithmetic step that forces an integral half-integer to have
actually integral coordinates.
\end{lemma}

\begin{lemma}[Carrier criterion for the Zsqrtd model]%
\label{lem:zsqrtd-carrier}
\lean{QuadraticNumberFields.RingOfIntegers.Zsqrtd.halfInt_mem_range_toQsqrtdHom_iff_even_even}\leanok
A half-integer element $(a' + b' \sqrt d)/2$ lies in the image of
$\mathbf{Z}[\sqrt d] \hookrightarrow Q(\sqrt d)$ if and only if $a'$ and $b'$ are
both even.
\end{lemma}

\begin{theorem}[Integral elements lie in the Zsqrtd model off the
\texorpdfstring{$1 \bmod 4$}{1 mod 4} branch]%
\label{thm:integral-in-zsqrtd}%
\lean{QuadraticNumberFields.RingOfIntegers.exists_zsqrtd_of_isIntegral_of_ne_one_mod_four}\leanok
\uses{thm:halfint-normal-form,lem:mod4-ne-one,lem:zsqrtd-carrier}
Assume $d$ is squarefree and $d \not\equiv 1 \pmod 4$. Then every element of
$Q(\sqrt d)$ integral over\/ $\mathbf{Z}$ already comes from the order
$\mathbf{Z}[\sqrt d]$.
\end{theorem}

\section{The \texorpdfstring{$1 \bmod 4$}{1 mod 4} Branch}

\begin{lemma}[Modulo-$4$ criterion in the $1 \bmod 4$ branch]%
\label{lem:mod4-one}
\lean{QuadraticNumberFields.RingOfIntegers.dvd_four_sub_sq_iff_same_parity_of_one_mod_four}\leanok
Assume $d$ is squarefree and $d \equiv 1 \pmod 4$. Then
\[
  4 \mid (a'^2 - d b'^2)
  \iff
  a' \equiv b' \pmod 2.
\]
Thus the integral elements are no longer forced into $\mathbf{Z}[\sqrt d]$; they
are allowed to have both coordinates odd.
\end{lemma}

\begin{lemma}[Carrier criterion for the Zomega model]%
\label{lem:zomega-carrier}
\lean{ZOnePlusSqrtOverTwo.halfInt_mem_carrierSet_iff_same_parity}\leanok
A half-integer element belongs to the image of
$\mathbf{Z}[(1+\sqrt{1+4k})/2] \hookrightarrow Q(\sqrt{1+4k})$ if and only if its
two numerator coordinates have the same parity.
\end{lemma}

\begin{theorem}[Integral elements lie in the Zomega model]%
\label{thm:integral-in-zomega}%
\lean{QuadraticNumberFields.RingOfIntegers.exists_zOnePlusSqrtOverTwo_of_isIntegral_of_one_mod_four}\leanok
\uses{thm:halfint-normal-form,lem:mod4-one,lem:zomega-carrier}
Let $d = 1 + 4k$ with $d$ squarefree and $d \neq 1$. Then every element of
$Q(\sqrt d)$ integral over\/ $\mathbf{Z}$ lies in the order
$\mathbf{Z}[(1+\sqrt d)/2]$.
\end{theorem}

\chapter{Ring of Integers Classification}

\begin{theorem}[Generic identification of the ring of integers]%
\label{thm:ring-of-integers-embedding}%
\lean{QuadraticNumberFields.RingOfIntegers.ringOfIntegers_equiv_of_embedding}\leanok
Suppose a commutative ring $R$ embeds into a number field $K$ and its image is
exactly the set of elements of $K$ integral over\/ $\mathbf{Z}$. Then $R$ is
canonically isomorphic to the ring of integers $\mathcal{O}_K$.
\end{theorem}

\begin{theorem}[The Zsqrtd order is maximal when
\texorpdfstring{$d \not\equiv 1 \pmod 4$}{d not congruent 1 mod 4}]%
\label{thm:roi-zsqrtd}%
\lean{QuadraticNumberFields.RingOfIntegers.ringOfIntegers_equiv_zsqrtd_of_mod_four_ne_one}\leanok
\uses{thm:integral-in-zsqrtd,thm:ring-of-integers-embedding,def:zsqrtd}
Assume $d$ is squarefree, $d \neq 1$, and $d \not\equiv 1 \pmod 4$. Then
\[
  \mathcal{O}_{Q(\sqrt d)} \cong \mathbf{Z}[\sqrt d].
\]
This is the first branch of the classical ring-of-integers classification.
\end{theorem}

\begin{theorem}[The Zomega order is maximal when
\texorpdfstring{$d \equiv 1 \pmod 4$}{d congruent 1 mod 4}]%
\label{thm:roi-zomega}%
\lean{QuadraticNumberFields.RingOfIntegers.ringOfIntegers_equiv_zOnePlusSqrtOverTwo_of_mod_four_eq_one}\leanok
\uses{thm:integral-in-zomega,thm:ring-of-integers-embedding,def:zomega}
Assume $d$ is squarefree, $d \neq 1$, and $d = 1 + 4k$. Then
\[
  \mathcal{O}_{Q(\sqrt d)} \cong \mathbf{Z}\Bigl[\frac{1+\sqrt d}{2}\Bigr].
\]
This is the second branch of the classification theorem.
\end{theorem}

\begin{theorem}[Classification of the ring of integers of quadratic fields]%
\label{thm:roi-classification}%
\lean{QuadraticNumberFields.RingOfIntegers.ringOfIntegers_classification}\leanok
\uses{thm:roi-zsqrtd,thm:roi-zomega}
Let $d$ be a squarefree integer with $d \neq 1$. Then exactly one of the
following descriptions of $\mathcal{O}_{Q(\sqrt d)}$ holds:
\begin{description}
\item[Non-$1 \bmod 4$ branch] If $d \not\equiv 1 \pmod 4$, then
  $\mathcal{O}_{Q(\sqrt d)} \cong \mathbf{Z}[\sqrt d]$.
\item[$1 \bmod 4$ branch] If $d \equiv 1 \pmod 4$, then there exists $k$ with
  $d = 1 + 4k$ and
  $\mathcal{O}_{Q(\sqrt d)} \cong \mathbf{Z}[(1+\sqrt d)/2]$.
\end{description}
This is the main theorem formalized by the repository.
\end{theorem}

\chapter{A Structural Corollary}

\begin{theorem}[Failure of Dedekind-ness in the $1 \bmod 4$ branch]%
\label{thm:not-dedekind-one}%
\lean{QuadraticNumberFields.RingOfIntegers.not_isDedekindDomain_zsqrtd_of_mod_four_eq_one}\leanok
\uses{def:zsqrtd,def:zomega}
If $d$ is squarefree, $d \neq 1$, and $d \equiv 1 \pmod 4$, then the order
$\mathbf{Z}[\sqrt d]$ is not Dedekind. Concretely, the integral element
$(1+\sqrt d)/2$ witnesses that $\mathbf{Z}[\sqrt d]$ is not integrally closed in
its fraction field.
\end{theorem}

\begin{corollary}[Dedekind criterion for the Zsqrtd order]%
\label{cor:dedekind-criterion}%
\lean{QuadraticNumberFields.RingOfIntegers.isDedekindDomain_zsqrtd_iff_mod_four_ne_one}\leanok
\uses{thm:roi-zsqrtd,thm:not-dedekind-one}
Let $d$ be a squarefree integer with $d \neq 1$. Then
\[
  \mathbf{Z}[\sqrt d] \text{ is a Dedekind domain }
  \iff
  d \not\equiv 1 \pmod 4.
\]
This packages the classification theorem into the exact criterion used later by
the examples and by the comparison with mathlib's \verb|Zsqrtd|.
\end{corollary}

\chapter{Module Map}

The code underlying this blueprint is concentrated in a small group of files:
\begin{description}
\item[\texttt{Basic.lean}] Defines \verb|Qsqrtd| and the ambient quadratic-field
model.
\item[\texttt{RingOfIntegers/TraceNorm.lean}] Shows that integral elements have
integral trace and norm, and that their imaginary part has denominator at most
$2$.
\item[\texttt{RingOfIntegers/HalfInt.lean}] Defines the half-integer notation and
computes its trace and norm.
\item[\texttt{RingOfIntegers/ModFour.lean}] Proves the arithmetic dichotomy
modulo $4$.
\item[\texttt{RingOfIntegers/Integrality.lean}] Converts the trace/norm data into
branch-specific membership theorems for the two integral models.
\item[\texttt{RingOfIntegers/Classification.lean}] Upgrades those membership
results to ring-of-integers isomorphisms and Dedekind consequences.
\end{description}

This leaves a clean downstream interface: once
\cref{thm:roi-classification,cor:dedekind-criterion} are available, the
discriminant formulas, the Dedekind-domain instances, and the
$\mathbf{Z}[\sqrt{-5}]$ examples can all work with the correct maximal order.
